
\documentclass[template=tabling,81pt,headonall]{azmoon}
\usepackage{xepersian}
\usepackage{amsfonts}
\usepackage{graphicx}
\usepackage{svg}
\svgpath{ {./images/} }
\graphicspath{ {./images/} }
\settextfont{Yas}
\setdigitfont{A Iranian Sans}
\usepackage{fontawesome5}

\printanswers
    \teacher{محمد صالح علی اکبری}
    \teachertitle{دبیر}
    \city{گناباد}
    \schooltitle{دبیرستان دوره اول}
    \school{شهید جوانبخت}
    \grade{ریاضی هفتم (نمونه ۲)}
    \branch{نمونه ۲}
    \topic{فصل ۱: راهبردهای حل مسئله}
    \examdate{۲۳ مهر ۱۴۰۵}
    \answertime{۴۰ دقیقه}
    \begin{document}
	\begin{questions}
		\nointerlineskip%
		\vskip-\baselineskip
		\question{%
در الگوی زیر، چوب‌کبریت‌ها برای ساختن مثلث‌های به هم چسبیده به کار رفته‌اند:

\noindent\textbf{شکل (۱):} $3$ چوب‌کبریت \hfill \textbf{شکل (۲):} $5$ چوب‌کبریت \hfill \textbf{شکل (۳):} $7$ چوب‌کبریت

الف) شکل چهارم با چند چوب‌کبریت ساخته می‌شود؟ \\
ب) رابطهٔ جبری بین شمارهٔ شکل ($n$) و تعداد چوب‌کبریت‌ها را بنویسید. \\
ج) شکل بیست و پنجم ($25$) از چند چوب‌کبریت تشکیل شده است؟‌
\\‌
\\‌
\\}\question{%
توپی از ارتفاع $18$ متری رها می‌شود. این توپ پس از هر بار برخورد به زمین، تا $\frac{2}{3}$ ارتفاع قبلی خود بالا می‌آید. این توپ پس از برخورد سوم به زمین، تا چه ارتفاعی بالا خواهد آمد؟ \textbf{(حل مسئله با «رسم شکل جهش‌ها» و نوشتن مراحل محاسبه الزامی است)}‌
\\‌
\\‌
\\‌
\\‌
\\}\question{%
در یک مزرعه در مجموع $20$ رأس مرغ و گوسفند نگهداری می‌شود. اگر تعداد کل پاهای آن‌ها روی هم $58$ عدد باشد، در این مزرعه چند مرغ و چند گوسفند وجود دارد؟ \textbf{(جدول حدس و آزمایش را با حداقل سه حدس رسم کرده و پاسخ نهایی را بنویسید)}‌
\\‌
\\‌
\\‌
\\‌
\\}\question{%
مجموع سن سارا و خواهرش برابر $16$ سال است. اگر حاصل‌ضرب سن آن‌ها بیشترین مقدار ممکن باشد، سن هر کدام چند سال است؟ \textbf{(جدول حالت‌های ممکن جمع را به صورت نظام‌دار بنویسید و حالت مطلوب را مشخص کنید)}‌
\\‌
\\‌
\\‌
\\}\end{questions}
    \end{document}
    